\section{Method}
\label{sec:method}

\paragraph{Overview.}
\method separates two decisions: which precision configuration fits the storage budget, and how that configuration quantizes inputs during execution. We first state the closed-loop evaluation objective, then describe the initializer and complete-policy deviation audit used for offline adaptation. Input-conditioned A8 supplies the deployment rule. The output is one frozen mask and one activation rule; neither success feedback nor task routing is available at runtime.

\subsection{The Full-Context Gap under an Exact Storage Budget}
\label{sec:problem}

Let $f_{\theta}(o,\epsilon)$ denote a frozen FP16-reference VLA that maps observation $o$ and generative noise $\epsilon$ to an action chunk $A\in\mathbb{R}^{H\times d}$. The target linear layers are indexed by $\mathcal{L}$. A mask $M\in\{4,16\}^{|\mathcal{L}|}$ assigns signed-nibble W4 or native 16-bit precision to each target, and $Q_M(f_\theta)$ denotes the mixed-precision policy under a fixed activation rule. Here FP16 follows the reference-row naming; native tensors may use BF16 in the model adapter. Let $\mathcal D_M$ be the rollout distribution induced by this policy under a fixed task distribution, and let $S(Q_M;\tau)\in\{0,1\}$ denote task success on rollout $\tau$. The deployment objective is
\begin{equation}
 M^\star=\arg\max_{M\in\{4,16\}^{|\mathcal L|}}
 \mathbb E_{\tau\sim\mathcal D_M}\!\left[S(Q_M;\tau)\right]
 \quad\text{s.t.}\quad C(M)\le C_{\max},
 \label{eq:problem}
\end{equation}
where $C(M)$ is exact packed static-component storage. Equation~\ref{eq:problem} states the evaluation goal, not the selection rule: success labels are unavailable during FCP, and collecting $\mathcal D_M$ for every mask is intractable. We therefore screen complete policies with label-free physical-sequence and component-deviation guards on teacher-induced contexts $\mathcal D_T$, then collect candidate-induced contexts only for shortlisted masks. Section~\ref{sec:fcp} constructs anchor $M_0$, defines the finite audited class $\mathcal C(M_0)$, and returns $\widehat M$ with abstention.

The dependence of both success and fidelity on $\mathcal D_M$ follows from standard closed-loop error propagation. For a local linearization with action error $e_t=a_t^Q-a_t^T$,
\begin{equation}
 \delta x_{t+1}=A_t\delta x_t+B_te_t,
 \qquad
 \delta x_T=\sum_{t=0}^{T-1}\Phi_{T,t+1}B_te_t,
 \label{eq:closed_loop_propagation}
\end{equation}
where $\Phi_{T,t+1}$ composes the intervening state Jacobians. Per-step MSE measures $\lVert e_t\rVert^2$ but ignores error propagation, directional coherence, and the observations reached after a deviation. This motivates a sequence-level control contract and candidate-state audit; closed-loop rollouts remain the final measure of task success.

FCP is label-free: success and reward never enter its guards. Task identifiers only stratify uncertainty; candidate and reference share generative noise; and an independent noise stream is reserved for post-freeze audit. The upstream initializer is not label-free because a disclosed four-task success sweep chooses one objective weight. The final mask and activation rule are global. No gradient update, task-conditioned routing, output correction, or FP16 bypass is permitted.

Before deployment, FCP may use the FP16 teacher and unlabeled on-policy observations; during deployment, it cannot observe success or change precision. The mask is thus a fixed, auditable artifact rather than a runtime selector.

\subsection{Action-Weighted CS--CKA Precision Allocation}
\label{sec:cscka}

The allocator asks where a compressed policy should spend its limited FP16 capacity. CKA measures the geometry retained by a layer intervention, whereas CS-Aligner measures its distribution shift~\cite{kornblith2019similarity,yin2026csaligner}. For target $i$, we combine jointly normalized geometry loss $\bar d_i^g=\operatorname{norm}(1-\widehat{\mathrm{CKA}}_i)$ and distribution shift $\bar d_i^d=\operatorname{norm}(\widehat D_{\mathrm{CS},i})$:
\begin{equation}
 \begin{aligned}
 S_i&=\frac{\lambda_g\bar d_i^g+\lambda_d\bar d_i^d}{\lambda_g+\lambda_d},\\
 M_{\mathrm{init}}&\approx\arg\min_{M}\sum_i \mathbf{1}[M_i=4]w_iS_i
 \quad\mathrm{s.t.}\quad C(M)\le C_{\max},\;M_i=16\;(i\in\mathcal G).
 \end{aligned}
 \label{eq:cscka_allocator}
\end{equation}
Here $w_i$ is a winsorized action-impact weight, $\mathbf{1}[M_i=4]$ charges the estimated cost of quantizing target $i$, and $\mathcal G$ contains targets excluded from W4 by RMS and saturation guards. Thus, precision follows the action-weighted joint cost of representation geometry and distribution stability under the actual byte constraint. Native allocation can already satisfy the ceiling, whereas transferring its protection pattern to a different architecture can violate the target budget. FCP handles this latter case by using complete-policy evidence to decide which protections to remove. A four-task development sweep fixes $\lambda_g:\lambda_d=16:1$ from seven choices. A greedy byte-constrained initialization followed by fixed-size flips approximates this objective (Appendix~\ref{sec:init_provenance}). This produces a starting configuration, whose necessity and signal fusion are tested rather than assumed.

We estimate the allocation evidence on FP16 on-policy observations from 12 tasks, with four tasks from each of Atomic-Seen, Composite-Seen, and Composite-Unseen. Three seeds per task and four replans per sequence yield 36 sequences and 144 observations. A task--seed hash assigns early, middle, or late rollout windows. Task identifiers stratify the statistics, but Equation~\ref{eq:cscka_allocator} returns one global mask.

\subsection{Full-Context Policy Audit}
\label{sec:dpac}

\paragraph{Physical-sequence contract.}
\dpac measures deviation from teacher actions for a chunked, replanning VLA policy. It compares ordered physical actions over a replan sequence through coherent bias, pose drift, replan discontinuities, gripper changes, and tail deviation. FCP uses this proxy to audit complete-policy interventions. It is neither a task-success estimator nor a guarantee of behavioral preservation.

\paragraph{Paired physical sequences.}
For each task--seed sequence, teacher and candidate share observations and initial action noise. We inverse-normalize the final chunks and retain all $H=16$ actions executed before each replan, yielding an ordered sequence without overlapping forecasts. Errors use one robust coordinate scale estimated globally from FP16 data, so a candidate cannot alter its denominator. Five terms then measure local action error, normalized prefix drift, right-composed SE(3) pose drift, replan-boundary discontinuity, and gripper state/event timing. The prefix term preserves coherent bias, and the last three expose control errors hidden by coordinatewise MSE. Appendix~\ref{sec:metric_specification} gives all formulas and edge cases.

The five normalized components receive equal weight:
\begin{equation}
 \begin{aligned}
 d_{\mathrm{PAC}}^{(u)}&=\tfrac{1}{5}\bigl(
 d_{\mathrm{local}}+d_{\mathrm{prefix}}+d_{\mathrm{pose}}+
 d_{\mathrm{stitch}}+d_{\mathrm{grip}}\bigr),\\
 \dpac&=\operatorname{mean}_u d_{\mathrm{PAC}}^{(u)}+
 \operatorname{CVaR}_{0.9,u}\!\left(d_{\mathrm{PAC}}^{(u)}\right).
 \end{aligned}
 \label{eq:dpac_sequence}
\end{equation}
Here $u$ indexes ordered task--seed sequences; the outer CVaR gives the worst 10\% the same total weight as the mean. We retain \dfunc as a complementary final-action and kinematic diagnostic, not a second closed-loop loss.

\paragraph{Abstaining complete-policy audit.}
\label{sec:fcp}

FCP treats mask adaptation as constrained decision-making rather than a second allocation heuristic. Complete-policy coordinate evidence retains or projects $M_{\mathrm{init}}$ to obtain a byte-feasible anchor $M_0$; paired evaluation of complete proposal masks then decides whether to admit a later refinement. Thus, abstention can still follow a consequential feasibility projection: it returns $M_0$, which equals $M_{\mathrm{init}}$ only when the initializer was already feasible.

\paragraph{Two-stage protocol.}
The construction stage pairs complete-policy sequences on a stratified buffer, evaluates instantiated one-layer interventions, projects an over-budget seed into the feasible set, and generates a finite set of exact-byte proposals. The acceptance stage applies task-level one-standard-error bounds to \dpac, \dfunc, pose, stitch, and gripper deviation. Any surviving proposal is checked again on candidate-induced observations before it can replace $M_0$.

\paragraph{Evidence-guided feasibility projection.}
If $C(M_{\mathrm{init}})\le C_{\max}$, FCP sets $M_0=M_{\mathrm{init}}$. Otherwise, it instantiates each FP16-to-W4 removal as a complete policy and removes protections from least to most damaging until the mask is feasible. This label-free projection ensures that abstention always returns a deployable mask; only the subsequent audit may replace $M_0$.

\paragraph{Coordinate counterfactuals.}
Starting from $M_0$, FCP flips target precision states one at a time and, when scheduled, repeats the coordinate sweep around the preceding winner. Each intervention is instantiated as a mixed-precision network and compared with its complete-mask background. The resulting conservative benefits feed an exact sparse Pareto 0--1 dynamic program subject to
\begin{equation}
 C(M)=C_{\mathrm{fixed}}+\sum_{i\in\mathcal{L}}C_i(M_i)
 \le 1.10\,C_{\mathrm{QuantVLA}}.
 \label{eq:budget}
\end{equation}
The constraint is expressed in bytes rather than layer count because layers differ in tensor and metadata size. The dynamic program only generates candidates: its additive counterfactual benefits cannot capture interactions among simultaneous layer changes.

The finite class $\mathcal C(M_0)$ contains $M_0$, dynamic-program proposals, and five structured proposal types. Their protected sets are rebuilt from the layer scores, not restricted to one-layer neighbors of $M_0$: the GR00T proposals retain only one to five FP16 layers versus 16 in $M_0$. A positive conservative FP16 benefit means that retaining or restoring FP16 is valuable; for an already protected layer it does not mean that removing protection improves the policy. Additive scores measured around $M_0$ may extrapolate poorly to these multi-layer changes. Complete-policy evaluation checks the emitted class, not the entire feasible neighborhood (Appendix~\ref{sec:fcp_visualization}).

\paragraph{Task-level minimax adjudication.}
For metric $q\in\{\dfunc,\dpac\}$ and task cluster $g\in\{\mathrm{all},\mathrm{Atomic},\mathrm{C\text{-}Seen},\mathrm{C\text{-}Unseen}\}$, let $\Delta_{q,t}(M)$ be the paired candidate-minus-$M_0$ task scalar. We define
\begin{equation}
 U_{q,g}(M)=
 \frac{\operatorname{mean}_{t\in g}\Delta_{q,t}(M)+
 \operatorname{SE}^{\mathrm{jack}}_{t\in g}[\Delta_{q,t}(M)]}
 {\operatorname{mean}_{t\in g}q_t(M_{\mathrm{base}})},
 \qquad
 J(M)=\max_{q,g}U_{q,g}(M).
 \label{eq:minimax}
\end{equation}
Here $M_{\mathrm{base}}=M_0$ is evaluated on the same context as the candidate. The jackknife makes the task, rather than an observation, the uncertainty unit. Eligibility requires $J(M)<0$ and non-positive one-SE upper bounds for pose, stitch, and gripper; these conservative guards are not physical-safety certificates.

\paragraph{Candidate-state check and freeze.}
The protocol reserves candidate-induced observations over four to eight replans for shortlisted candidates and defines $J_{\mathrm{final}}=\max(J_{\mathrm{teacher}},J_{\mathrm{candidate}})$. In the implemented selection order, eligible teacher-state candidates are first filtered by a paired one-standard-error comparison against the best teacher-state objective; up to three then enter candidate-state adjudication. Let $\mathcal E(M_0)$ denote the remaining byte-feasible candidates with negative $J_{\mathrm{final}}$ and the required deviation guards. The frozen choice is
\begin{equation}
\widehat M=\begin{cases}
\operatorname*{lexmin}_{M\in\mathcal E(M_0)}
 (C(M),N_{16}(M),\operatorname{names}(M)),&\mathcal E(M_0)\ne\varnothing,\\
M_0,&\mathcal E(M_0)=\varnothing,
\end{cases}
\label{eq:restricted_freeze}
\end{equation}
where $N_{16}$ counts protected layers and names provide deterministic ordering. Storage breaks the tradeoff only after deviation eligibility: the rule does not exchange increased proxy error for fewer bytes. On GR00T, all proposals fail the teacher-state screen, so candidate-state adjudication is not exercised. Retaining $M_0$ certifies neither local optimality nor prevention of a task-success regression.

\subsection{Deployment Context: Input-Conditioned Activation Quantization}
\label{sec:deployment}

The deployment stage removes the assumption that a frozen range table covers later closed-loop inputs. Hessian-aware W4 provides the weight substrate, while A8 uses a standard per-forward max-abs estimator. This simple rule isolates the deployment policy---conditioning ranges on realized inputs---rather than proposing a new scalar estimator.

\paragraph{Weight quantization.}
For each W4 layer, 256 calibration observations form FP16 input second-order statistics $X^\top X$. A GPTQ-style sequential error-feedback pass selects clipping and quantizes signed weights in groups of 64~\cite{frantar2022gptq}. Two signed four-bit codes are packed per byte; per-group scales and packing metadata are retained, but the FP16 weight tensor is not. Row rotation is the identity, so the reported method does not rely on an unreported learned transform.

\paragraph{Input-conditioned activation range.}
A calibrated static activation table can become mismatched as the observation, flow step, and policy-induced state change. Our static counterfactual uses independently calibrated 99.9th-percentile scales for each target layer, input channel, and applicable flow step. \dyrange removes this lookup and recomputes a scale for every target linear layer, forward call, and input channel:
\begin{equation}
 s_c=\max\!\left(\frac{\max |x_{\ldots,c}|}{127},10^{-6}\right),
 \qquad q_c=\operatorname{clip}\!\left(\operatorname{round}(x_c/s_c),-127,127\right).
 \label{eq:dyrange}
\end{equation}
The rule is deterministic and uses the input dtype. It has no calibration lookup, flow-step selector, task selector, or candidate branch. Equation~\ref{eq:dyrange} is applied uniformly, and runtime output correction is disabled. Appendix~\ref{sec:hardware_results} reports hardware measurements under explicitly audited runtime contracts.

\paragraph{Exact storage scope.}
We report the byte sum of fixed model components, packed W4 codes, per-group scales and metadata, and native FP16 tensors. This scope excludes activations, CUDA workspaces, simulator state, and eager fake-quant caches. The projected $\pi_{0.5}$ anchor contains 121 W4 and 59 FP16 targets and occupies 1,634,828,288 bytes; the native GR00T plan contains 100 W4 and 16 FP16 targets and occupies 962,068,480 bytes. These values establish static compression only. A separate A40 experiment measures runtime quantities directly.
